\section{Teaching and Academic Research}
\label{sec:teaching}

Across these institutions, used interactive, peer-led, and just-in-time teaching methods.

\record{Lorain County Community College}{Aug. 2016 - May 2019}{\textbf{Adjunct Professor} - Taught the laboratory component of algebra-based General Physics I (then PHYC 151) and the three-credit ASTY 152 Solar System Astronomy lecture course.}
\begin{itemize}
  \item Guided students through laboratory work in mechanics, energy, momentum, rotation, and thermal physics, emphasizing measurement uncertainty, quantitative analysis, and scientific reporting.
  \item Developed and taught online astronomy instruction covering astronomical observation, planetary-system formation, terrestrial and Jovian planets, small bodies, and space exploration.
  \item Exceeded institutional student-evaluation benchmarks by an average of 16\% in 2017 and 17\% in 2018.
  \item Reduced reported student confusion by 28\% by reorganizing online instructional materials and reducing the need for clarification during office hours.
  \item Developed and published more than 70 hours of publicly available recorded lectures and laboratory content to support flexible, on-demand learning.
  \item Used Just-in-Time Teaching strategies to adapt instruction for underperforming students based on pre-class responses.
\end{itemize}
\record{Cincinnati State Technical and Community College}{Feb. 2011 - Jun. 2011}{\textbf{Adjunct Professor} - Developed and delivered the integrated Technical Physics I--III lecture and laboratory sequence for technical-program students with varied mathematics preparation.}
\begin{itemize}
  \item Delivered two weekly lecture hours and three laboratory hours per course across applied mechanics, thermodynamics, electricity and circuits, waves, and optics.
  \item Wrote laboratory instructions and coordinated team-based measurement, data-analysis, and technical-report work.
\end{itemize}
\record{Miami University}{Aug. 2010 - Jun. 2011}{\textbf{Adjunct Professor} - Developed and taught three-credit general-education courses PHY 101 Physics and Society and PHY 111 Astronomy and Space Physics.}
\begin{itemize}
  \item Designed a concept-first applied-physics curriculum for non-science majors, translating advanced technical subjects into accessible, decision-ready frameworks without relying on calculus.
  \item Taught students to use order-of-magnitude estimates, energy density, power, conservation laws, and comparative risk to evaluate claims involving transportation fuels, electricity generation, climate, and emerging technologies.
  \item Connected nuclear physics and radioactivity with power generation, medical imaging, radiation risk, and nuclear security; linked gravity, waves, and the electromagnetic spectrum with satellites, GPS, radar, communications, and remote sensing.
  \item Structured lessons around consequential real-world questions, teaching students to distinguish scientific constraints from policy and value judgments and to communicate evidence-based conclusions to nontechnical audiences.
  \item Taught astronomy topics spanning planetary systems, stellar evolution, galaxies, cosmology, and the scientific search for extraterrestrial life.
\end{itemize}
\record{University of Cincinnati}{Jan. 2006 - Sep. 2010}{Held concurrent research and teaching appointments in physics, astronomy, and engineering computation.}
\subrole{Research Assistant, Department of Physics}{Sep. 2006 - Sep. 2010}
Conducted near-infrared astronomical research using observations from the 0.8--5.4 micrometer SpeX spectrograph at the NASA Infrared Telescope Facility.
\begin{itemize}
  \item Reduced and analyzed $0.8$--$2.4$ micrometer, $R \sim 2000$ SpeX spectra and measured diagnostic H~I, He~I, He~II, C~IV, N~III, and CO spectral features across nine candidate cluster members.
  \item Co-developed an astrometrically calibrated catalog of 130 sources by registering SpeX acquisition images to 2MASS; compiled available $JHK$ photometry and used DAOPHOT to separate blended sources in the crowded field.
  \item Helped classify the cluster's massive stellar population, identifying an O5~V dwarf, several early-B stars, and a massive young stellar object exhibiting CO band-head emission from dense circumstellar material.
  \item Combined spectral classifications with 2MASS photometry to estimate cluster extinction ($A_V \sim 21$--$23$), distance ($1.0$--$1.75$ kpc), age of no more than a few million years, and mass of approximately $10^3$ solar masses.
  \item Cross-matched infrared and radio-source positions and identified candidate near-infrared counterparts for three deeply embedded radio sources associated with early star formation.
  \item Published the resulting 130-source astrometric and photometric catalog through the international VizieR astronomical data service.
  \item Acquired follow-up L-band spectra of Galactic OB stars with SpeX during three consecutive nights at NASA's IRTF.
  \item Estimated stellar radial velocities by measuring pixel-space shifts of CO absorption features against normalized standard-star spectra.
  \item Performed astronomical spectroscopy, photometry, classification, calibration, signal analysis, and noise reduction using IRAF, IDL, SpeXtool, Mathematica, and Perl.
\end{itemize}
\begin{minipage}{\linewidth}
\subrole{Teaching Assistant}{Jan. 2006 - Jun. 2010}
Led approximately 60 laboratory meetings for more than 100 engineering, science, and preprofessional students in the algebra/trigonometry-based College Physics sequence (PHYS 111--113) and calculus-based General Physics sequence (PHYS 211--213).
\end{minipage}
\begin{itemize}
  \item Taught experimental design, measurement uncertainty, data acquisition, graphical analysis, and technical reporting using PASCO DataStudio and Microsoft Excel.
  \item Served as teaching assistant and grader for the corresponding College Physics lecture sequence (PHYS 101--103) and General Physics lecture sequence (PHYS 201--203).
  \item Supported Astronomy and the Natural Universe lectures (PHYS 121--122) and astronomy laboratories (PHYS 125--127) through grading and instructional assistance.
  \item Designed interactive laboratory instruction using Peer Instruction and Interactive Lecture Demonstrations to identify and address student misconceptions in real time.
  \item Supported group recitations in freshman engineering physics through demonstrations, short lectures, and individual assistance.
\end{itemize}
\subrole{Engineering Teaching Assistant}{Oct. 2009 - Dec. 2009}
Prepared and delivered Mathematica instruction for Numerical Methods for Engineering Design, an upper-division aerospace-engineering course.
\begin{itemize}
  \item Developed step-by-step demonstrations applying symbolic and numerical computing to interpolation and approximation, numerical integration, nonlinear equation solving, systems of equations, and ordinary differential equations.
  \item Connected computational methods with aerospace analysis and design problems, showing students how to translate mathematical formulations into executable workflows and interpret the resulting output.
  \item Reinforced numerical accuracy, error, stability, and method selection so students could evaluate the credibility of computed engineering results.
\end{itemize}
\record{Ohio University}{Oct. 2003 - May 2004}{\textbf{Student Researcher, Fine Arts} - Earned a competitive Provost's Undergraduate Research Fund award to adapt a wintergreen-oil toner-transfer process from lithography to intaglio photo-etch printing.}
\begin{itemize}
  \item Developed and tested a complete copper- and zinc-plate workflow spanning surface preparation, asphaltum grounding, high-contrast toner transfer, burnishing, clearing, aquatint, and staged etching.
  \item Replaced xylene- and acetone-based photo-transfer steps with methyl salicylate, eliminated the need for specialized exposure equipment, and demonstrated a process using less than \$20 in materials.
  \item Characterized critical process variables and failure modes, including ground thickness and cure time, solvent saturation, image resolution, transfer pressure, evaporation time, and temperature control during aquatint fusion.
\end{itemize}
